\section{Formal and statistical detail}
\label{app:formal}

\subsection{Notation and the assumption-bound relationship between conditional information and unique variance}
\label{sec:cmi-detail}

Let $Y$ denote a scalar brain-response target, $X$ the language-model feature block, and $Z$ the nuisance feature block.
Conditional mutual information asks how much statistical dependence remains between $X$ and $Y$ after $Z$ is known:
\begin{equation}
I(X;Y\mid Z)=H(Y\mid Z)-H(Y\mid X,Z).
\end{equation}
This definition is distributional and does not specify a linear estimator.
Under a population multivariate Gaussian model with correctly specified unregularized conditional means, the reduction in conditional residual variance can be written as
\begin{equation}
I(X;Y\mid Z)=-\frac{1}{2}\log\left(1-R^2_{\mathrm{partial}}\right).
\end{equation}
Here $R^2_{\mathrm{partial}}$ measures the fraction of the variance left after $Z$ that is explained by adding $X$.

The empirical score used in this thesis is instead semipartial cross-validated unique variance,
\begin{equation}
\uniqueR=R^2_{\mathrm{full}}-R^2_{\mathrm{nuisance}}.
\end{equation}
Under the same population assumptions, partial and semipartial coefficients are related by
\begin{equation}
\uniqueR=(1-R^2_{\mathrm{nuisance}})R^2_{\mathrm{partial}}.
\end{equation}
The thesis fits regularized ridge models on finite folds and evaluates them out of sample, so neither equality makes the measured $\uniqueR$ an estimator or lower bound for conditional mutual information.
The information-theoretic formulation supplies a conceptual target, while the actual estimand is incremental held-out linear prediction beyond the specified nuisance design.

The nuisance set is also necessarily incomplete.
Length, position, static embeddings, acoustic tiers, and temporal controls block named alternative pathways, but they do not condition on every aspect of stimulus content.
A positive $\uniqueR$ therefore licenses the statement that contextual features add held-out linear predictive value beyond those controls.
It does not isolate a neural computation, establish causal equivalence, or distinguish prediction of a shared high-level stimulus component from prediction of participant-specific processing.

\subsection{Regularization, data processing, generalization, and the rate--distortion analogy}

An auxiliary brain term gives an optimization objective of the form
\begin{equation}
\widehat\theta=\arg\min_\theta\left\{\mathcal L_{\mathrm{text}}(\theta)+\lambda\mathcal L_{\mathrm{brain}}(\theta)\right\}.
\end{equation}
The brain term biases optimization among solutions that also fit the text objective.
It can be read as a maximum-a-posteriori prior only if $\exp[-\lambda\mathcal L_{\mathrm{brain}}(\theta)]$ is normalizable and independent of the observed text likelihood.
Without those conditions it is more precise to call it an auxiliary loss or regularizer.

The data-processing inequality states that $I(A;C)\le I(A;B)$ for a justified Markov chain $A\rightarrow B\rightarrow C$.
If a compressed representation $C$ is literally a fixed post-processing of teacher representation $B$, then its information about a brain variable $A$ cannot exceed the teacher's.
Ordinary knowledge distillation does not automatically instantiate this chain because both teacher and student receive direct stimulus or corpus information.
The inequality therefore motivates a preservation question but does not predict the sign of the empirical encoding-score change under training.

For a learning algorithm mapping a sample $Z^n$ to parameters $W$, a sub-Gaussian information-theoretic generalization bound has the form
\begin{equation}
\left|\mathbb E[L(W)-\widehat L(W)]\right|\le \sqrt{\frac{2\sigma^2 I(W;Z^n)}{n}}.
\end{equation}
The right-hand side grows with algorithm--sample mutual information.
The bound does not imply that adding a neural signal creates a benefit, that the benefit scales as an added information budget, or that a low-data regime must improve.
An auxiliary constraint could tighten this bound only if it reduces dependence on the sampled text while preserving empirical fit, which was not measured here.
Earlier order-of-magnitude information-budget calculations are therefore not part of the thesis argument.

Classical rate--distortion theory defines
\begin{equation}
R(D)=\min_{P_{\widehat S\mid S}:\mathbb E[d(S,\widehat S)]\le D}I(S;\widehat S).
\end{equation}
The empirical experiments do not identify this object because parameter count or FLOPs are not the mutual-information rate and alignment loss is not established as the source-coding distortion.
The phrase ``rate--distortion-style trade-off'' is used only for the engineering Pareto relation among model budget, language-model quality, and measured alignment.

\subsection{Inference units and power}
\label{sec:inference-detail}

An inference unit is the independently varying entity over which an uncertainty statement generalizes.
Training seeds quantify optimizer variability, contiguous test folds quantify sensitivity to stimulus blocks, voxels quantify spatial measurement breadth, and participants quantify biological population variation.
These units answer different questions and cannot be substituted by multiplying their counts.

E026 uses two disjoint deterministic item samples to test the stability of fixed target-geometry summaries; they are not population replicates and receive no confidence intervals or $p$ values \evd{E026}.
Its six paired seeds are technical realizations of the frozen training procedure.
Reported seed intervals are descriptive $t$ intervals, and the exact $2^6$ magnitude-preserving sign-flip test relies on exchangeability of paired seed-contrast signs under the sharp null; neither device licenses biological or target-population inference.

E004 and E005 used three seeds and five rotating folds, producing 15 observed seed--fold cells.
Because seeds within one fold share the same held-out stimulus block and fold effects dominate the contrast, treating all 15 cells as independent understates uncertainty.
E004's mean Qwen contrast was \result{e004_qwen_specific}, but the five-fold $t$ interval was \result{e004_qwen_fold_ci}, and leaving out the influential fourth fold reduced the mean to \result{e004_qwen_lofo}. \evd{E004}
The correct conclusion is no demonstrated fold-level intervention effect, with limited power to exclude a small effect.

E008 instead took the predeclared median over the 15 paired seed--fold contrasts within each of nine participants and treated participants as the biological inference unit.
Its mean was \result{e008_mean} with participant-axis $t$ interval \result{e008_ci}. \evd{E008}
A $+0.003$ reference effect lay well outside that interval, and the post-run calculation reported power 1.0 at that effect size while leaving smaller effects possible. \evd{E008}

E006 illustrates a distinct error.
Bootstrapping 11,442 spatially dependent voxels generated a narrow interval but did not create participant replication.
The retained evidence is the conditional UTS03 point gap, \result{e006_gpt2_gap} for GPT-2 and \result{e006_qwen_gap} for Qwen, plus positive signs on all five held-out story blocks. \evd{E006}
Because those blocks share training stories, even the five signs are sensitivity evidence rather than an independent population confidence interval.

Power calculations must use the variance of the future paired contrast at its intended inference unit.
The retracted E007 calculation used unmatched absolute fold summaries from E006 and therefore did not estimate the variance of a future tuned-minus-control effect.
No replacement participant count is inferred from it.

\subsection{E005 reanalysis and the lambda sweep}
\label{sec:e005-detail}

E005 compared real-target brain-guided KD with a block-permuted-target twin while matching the text-training setup and approximately matching perplexity.
The initial averaged-target estimate was \result{e005_apparent}. \evd{E005}
Its original bootstrap sampled 15 seed--fold cells independently and excluded zero, but the corrected fold-unit analysis showed that one held-out block carried a disproportionate share of the result and the valid interval included zero.
The retained description is therefore a small averaged-target trend rather than a confirmed brain-specific improvement.

A later Qwen lambda sweep asked whether stronger auxiliary weighting produced a stable alignment gain without unacceptable language-model degradation.
Larger weights moved the training endpoint more strongly but also worsened perplexity, and no setting established a free target-specific improvement at matched quality. \evd{E005}
This supports a trade-off description rather than a favorable frontier shift.
The sweep is contextual because it reuses the participant-averaged target and does not replace E008's participant-level test.

\subsection{E015 family, band, and partial-correlation sensitivities}
\label{sec:e015-detail}

E015 scored 22 decoder language models from six named families using tokenizer-independent bits per byte and the same controlled Tuckute alignment endpoint.
Across the full capability range, the Pearson relationship was approximately $r=-0.78$. \evd{E015}
Removing near-floor models weakened the association: in the capable band it was approximately $-0.50$, and in the overlapping quality band approximately $-0.48$. \evd{E015}
The Fisher intervals for these band-restricted correlations included zero because the bands contained only eight to ten models.

A partial correlation controlling for log parameter count was $-0.675$ with Fisher 95\% interval $[-0.857,-0.344]$. \evd{E015}
This is not a clean decomposition into scale and quality because bits per byte and scale were strongly collinear and the residual variation also encoded family, data, and architecture differences.
Family-cluster resampling remained negative, but the six families separated into roughly two model generations, limiting the effective independent breadth.
Participant-level analyses showed the pattern in the one participant with adequate alignment signal and were uninformative in two participants near the measurement floor.
The supported conclusion is that quality is a serious confound requiring matched-quality controls, not that perplexity deterministically explains brain alignment within the capable regime.

\subsection{E017 intervention sequence and E023 identification gate}
\label{sec:e017-detail}

E017 tested whether full fine-tuning could reveal movement that low-rank adaptation missed.
An aggressive configuration substantially worsened held-out language quality and failed the quality requirement.
A gentler UTS03 pilot showed that the representation could move without the same failure, which cleared a three-participant existence test.
Across UTS01, UTS02, and UTS03, the real-target arm did not improve over its target-specific controls. \evd{E017}
The observed subject--seed cells share only three biological units, so the result is reported descriptively as a three-participant mechanism probe rather than with a population confidence claim.

E023 then asked how an objective-specific loss could be separated from the alignment expected at the student's quality.
Its proposed estimand was the vertical residual of a KD student below a same-lineage alignment--quality frontier, paired against ordinary fine-tuning from the same initialization.
The analysis-only slope gate found one locally flat lower-quality Pythia pair, but no Qwen pair met the predeclared saturated-regime requirement \evd{E023}.
Because the identifying region was not established, the training comparison was not launched.
Together E017 and E023 distinguish a tested full-fine-tuning null from an untested larger-scale counterfactual program.

\subsection{Compact notation}

\begin{description}[style=nextline,leftmargin=2.5cm] % notation layout, not a scientific result; \evd{E002}
\item[$S$] Language stimulus, such as a sentence or story segment.
\item[$X$] Language-model feature block for the stimulus.
\item[$Y$] Recorded or synthetic target response.
\item[$Z$] Specified nuisance feature block.
\item[$\uniqueR$] Cross-validated semipartial unique $R^2$, the full-model $R^2$ minus nuisance-model $R^2$.
\item[$\Lbrain$] Auxiliary loss relating a student representation to a neural or cognitive target.
\item[$\Eys$] Stimulus-predictable conditional mean of a response.
\item[ROI] Region of interest, a prespecified group or aggregate of voxels.
\item[BPB] Bits per byte, a tokenizer-independent language-model quality measure on a fixed corpus.
\item[KD] Knowledge distillation, training a student to match a teacher's softened output distribution.
\end{description}
